%% 
%% Copyright 2019-2024 Elsevier Ltd
%% 
%% This file is part of the 'CAS Bundle'.
%% --------------------------------------
%% 
%% It may be distributed under the conditions of the LaTeX Project Public
%% License, either version 1.3c of this license or (at your option) any
%% later version.  The latest version of this license is in
%%    http://www.latex-project.org/lppl.txt
%% and version 1.3c or later is part of all distributions of LaTeX
%% version 1999/12/01 or later.
%% 
%% The list of all files belonging to the 'CAS Bundle' is
%% given in the file `manifest.txt'.
%% 
%% Template article for cas-dc documentclass for 
%% double column output.

\documentclass[a4paper,fleqn]{cas-dc}

% If the frontmatter runs over more than one page
% use the longmktitle option.

%\documentclass[a4paper,fleqn,longmktitle]{cas-dc}

%\usepackage[numbers]{natbib}
%\usepackage[authoryear]{natbib}
\usepackage[authoryear,longnamesfirst]{natbib}

%%%Author macros
\def\tsc#1{\csdef{#1}{\textsc{\lowercase{#1}}\xspace}}
\tsc{WGM}
\tsc{QE}
%%%

% Uncomment and use as if needed
%\newtheorem{theorem}{Theorem}
%\newtheorem{lemma}[theorem]{Lemma}
%\newdefinition{rmk}{Remark}
%\newproof{pf}{Proof}
%\newproof{pot}{Proof of Theorem \ref{thm}}

\begin{document}
\let\WriteBookmarks\relax
\def\floatpagepagefraction{1}
\def\textpagefraction{.001}

% Short title
\shorttitle{}    

% Short author
\shortauthors{}  

% Main title of the paper
\title [mode = title]{RouteSAM: Semantic Routing for Semi-Supervised 2D Medical Image Segmentation via Cross-Image Propagation}  


% Title footnote mark
% eg: \tnotemark[1]
\tnotemark[1] 

% Title footnote 1.
% eg: \tnotetext[1]{Title footnote text}
\tnotetext[1]{} 

% First author
%
% Options: Use if required
% eg: \author[1,3]{Author Name}[type=editor,
%       style=chinese,
%       auid=000,
%       bioid=1,
%       prefix=Sir,
%       orcid=0000-0000-0000-0000,
%       facebook=<facebook id>,
%       twitter=<twitter id>,
%       linkedin=<linkedin id>,
%       gplus=<gplus id>]

\author[1]{}%[<options>]

% Corresponding author indication
\cormark[1]

% Footnote of the first author
\fnmark[1]

% Email id of the first author
\ead{}

% URL of the first author
\ead[url]{}

% Credit authorship
% eg: \credit{Conceptualization of this study, Methodology, Software}
\credit{}

% Address/affiliation
\affiliation[1]{organization={},
            addressline={}, 
            city={},
%          citysep={}, % Uncomment if no comma needed between city and postcode
            postcode={}, 
            state={},
            country={}}

\author[2]{}%[]

% Footnote of the second author
\fnmark[2]

% Email id of the second author
\ead{}

% URL of the second author
\ead[url]{}

% Credit authorship
\credit{}

% Address/affiliation
\affiliation[2]{organization={},
            addressline={}, 
            city={},
%          citysep={}, % Uncomment if no comma needed between city and postcode
            postcode={}, 
            state={},
            country={}}

% Corresponding author text
\cortext[1]{Corresponding author}

% Footnote text
\fntext[1]{}

% For a title note without a number/mark
%\nonumnote{}

% Here goes the abstract
\begin{abstract} 
The high cost of medical image annotation has made semi-supervised learning an important paradigm for medical image segmentation. However, when limited labeled data fail to sufficiently cover diverse lesion representations, semi-supervised segmentation models often suffer substantial performance degradation. This problem becomes particularly pronounced when large semantic discrepancies exist between labeled and unlabeled images, making it difficult to reliably transfer limited supervision and consequently reducing the quality of the resulting pseudo-supervision.
To address this issue, we propose RouteSAM, a progressive supervision propagation method based on semantic routing. To maximize feature-space coverage under limited supervision, RouteSAM first selects samples with high feature coverage as reliable starting points for subsequent propagation. For each unlabeled image, RouteSAM then constructs a propagation route composed of feature-related images between the supervised source and the target image. Supervision is progressively transferred through semantically neighboring images along the route, avoiding difficult direct transfer between images with large feature discrepancies.
Building upon this routing mechanism, we further introduce a multi-stage generalist–specialist collaboration strategy that exploits the complementary strengths of the general-purpose segmentation capability of SAM3 and the task-specific knowledge learned by a U-Net. Through iterative pseudo-supervision filtering and collaborative model adaptation, the proposed strategy progressively improves supervision reliability and enhances adaptation to the target medical segmentation task.
Experiments on Kvasir-SEG, BUSI, and ISIC2018 demonstrate the effectiveness of RouteSAM for semi-supervised medical image segmentation under extremely limited supervision, [final quantitative results to be inserted].
\end{abstract}

% Use if graphical abstract is present
%\begin{graphicalabstract}
%\includegraphics{}
%\end{graphicalabstract}

% Research highlights
\begin{highlights}
\item We propose a cross-image semantic routing and supervision propagation mechanism for semi-supervised 2D medical image segmentation. Instead of treating unlabeled samples as independent training instances, RouteSAM explicitly exploits cross-sample semantic relationships within the training set and constructs an ordered semantic route for each queried target, consisting of a labeled anchor, several semantically related unlabeled bridge images, and the target itself. SAM3 progressively transfers segmentation information along these routes, thereby expanding sparse manual annotations into large-scale cross-image pseudo supervision.

\item We introduce a coverage-aware propagation source selection strategy to improve semantic routing under extremely limited annotation budgets. In RouteSAM, labeled samples serve not only as direct supervision but also as starting points for cross-image propagation. We therefore select samples with broad coverage of the training distribution for annotation, providing more representative propagation sources across diverse lesion appearances and data variations.

\item We develop a propagation-oriented multi-stage generalist–specialist collaborative learning strategy between SAM3 and a task-specific U-Net. Route-derived pseudo supervision is first used to train the specialist, whose medical-domain predictions are then employed to audit the propagated supervision and guide SAM3 adaptation. The adapted SAM3 subsequently re-propagates supervision along the same semantic routes, producing refined pseudo labels for further specialist learning. This process progressively absorbs cross-image propagation knowledge into the segmentation models, enabling direct single-image inference without reference retrieval or semantic route construction at test time. 
\end{highlights}


%\nocite{*}

% Keywords
% Each keyword is seperated by \sep
\begin{keywords}
 \sep \sep \sep
\end{keywords}

\maketitle

% Main text



\section{Introduction}

Medical image segmentation provides fine-grained spatial delineation of lesions, organs, and other anatomically meaningful structures, and serves as an important foundation for computer-aided diagnosis, quantitative disease assessment, and treatment planning. In recent years, deep learning has significantly improved the performance of medical image segmentation, but high-performance segmentation models usually require a large amount of high-quality pixel-level annotations. Compared with natural images, accurate annotation of medical images often requires clinicians to invest substantial time and possess domain-specific expertise, making large-scale medical image annotation costly and difficult to obtain~\cite{hesamian2019deep,tajbakhsh2020embracing}. To reduce the dependence on manual annotation, semi-supervised medical image segmentation (SSMIS) has become an important research direction by jointly exploiting a small amount of labeled data and a large amount of unlabeled data~\cite{jiao2024learning}.

Existing SSMIS methods mainly obtain additional supervision from unlabeled data through pseudo-label learning and consistency regularization. One line of methods constructs pseudo-labels from model predictions and continuously updates the supervision information via teacher--student architectures or collaborative learning among multiple networks. Another line of methods perturbs the input, feature representations, or the model itself, and enforces prediction consistency under different perturbations. Representative methods such as Cross Teaching, MC-Net+, and BCP improve the utilization of unlabeled data through mutual supervision between heterogeneous networks, uncertainty-aware consistency learning, and bidirectional mixing between labeled and unlabeled samples, respectively~\cite{luo2022crossteaching,wu2022mcnet,bai2023bcp}. However, when the amount of manually labeled samples further decreases, the reliable task knowledge that the model can learn from labeled data also declines accordingly, making the pseudo-labels generated by the model itself more susceptible to prediction errors and uncertainty.

The emergence of promptable segmentation foundation models has brought new opportunities for obtaining stronger supervision. The Segment Anything Model (SAM) and its subsequent variants acquire general visual representation capabilities and prompt-driven segmentation priors through large-scale pre-training, and can segment target regions based on prompts such as points, boxes, or masks~\cite{kirillov2023segment}. Although the domain gap between natural images and medical images makes the direct segmentation performance of SAM unstable across different medical imaging tasks~\cite{mazurowski2023samexperiment,ma2024medsam}, the general segmentation knowledge obtained through large-scale pre-training remains highly valuable. Therefore, recent studies have begun to introduce SAM into semi-supervised medical image segmentation: task-specific models are used to generate prompts, and the predictions of SAM are used to construct or improve pseudo-supervision on unlabeled data. For example, SemiSAM uses a task-specific segmentation model to provide localization prompts for SAM; SemiSAM+ further establishes a collaborative relationship between a general foundation model and a task-specific specialist model; and SAMatch leverages SAM to further refine existing predictions~\cite{zhang2023semisam,zhang2025semisamplus,xu2024samatch}. These studies suggest that the general segmentation priors contained in foundation models can serve as another source of supervision for learning from unlabeled medical images. However, most existing methods still focus on improving the prediction or pseudo-label quality of each individual unlabeled image, and have not fully exploited the cross-sample semantic relationships among training samples.

It is worth noting that medical images are not independent of each other. The same anatomical structure or similar lesions often exhibit correlated semantic patterns across different samples. Existing studies have shown that cross-sample semantic relationships among different medical images can be exploited. For example, ACTION simultaneously models global semantic relationships and local anatomical information to improve semi-supervised representation learning, indicating that training samples do not necessarily need to be treated as completely independent instances~\cite{you2022action}. Furthermore, YoloSeg uses SAM2 to propagate segmentation information from a single labeled image to unlabeled images, demonstrating the potential of cross-image propagation in low-annotation regimes~\cite{zhang2026yoloseg}. However, such propagation is still mainly performed in a single-hop manner, directly from a labeled image to the target image. When two images differ substantially in lesion morphology, scale, or appearance, direct propagation needs to bridge a large inter-image gap in one step, which can easily lead to unstable pseudo-supervision. This raises a further question: can we exploit the sample relationships within unlabeled data to insert several semantically related intermediate images between the existing supervision and a difficult target, thereby decomposing one difficult direct propagation into several more local propagation steps?

To address this question, we further study how to organize the propagation of limited supervision by exploiting the relationships within unlabeled data. We view the unlabeled training set as a latent space formed by semantic relationships among samples, and propose \textbf{RouteSeg}. Under a fixed annotation budget, RouteSeg first selects labeled samples that can broadly cover the training data distribution, and defines them as \textbf{Anchors}. For each unlabeled \textbf{Target} whose pseudo-label needs to be generated, RouteSeg retrieves several semantically related images from the unlabeled data as \textbf{Bridges}, and arranges them into an ordered \emph{semantic route} according to their cross-sample relationships:

\begin{equation}
A \rightarrow B_1 \rightarrow B_2 \rightarrow \cdots \rightarrow B_K \rightarrow T.
\end{equation}

Then, RouteSeg performs \emph{progressive cross-image propagation} along this semantic route based on SAM3, gradually propagating the segmentation information provided by the Anchor to the Target. Compared with direct propagation, the introduction of Bridges decomposes a large inter-image variation into multiple more local propagation processes. In this way, an unlabeled sample is no longer merely a Target waiting for pseudo-label generation; it can also serve as a Bridge in the semantic route of another Target, thereby participating in the expansion of limited supervision. Since different Anchors and semantic routes may produce different mask candidates for the same Target, directly using all propagation results may introduce unreliable pseudo-supervision. To address this, RouteSeg further evaluates the quality of these mask candidates and selects reliable \emph{route-derived pseudo-labels} for subsequent training. The selected pseudo-labels and the ground-truth annotations of the Anchors are jointly used to train a task-specific segmentation model, which is referred to as the \textbf{specialist}; SAM3 serves as the \textbf{generalist}. The specialist learns task-specific knowledge from the target medical data and further provides feedback to the propagation results of the generalist, allowing the general segmentation priors of the foundation model to complement the domain knowledge of the medical task. Semantic routing and cross-image propagation are used only during training; at inference time, the specialist can directly segment a single image without retrieving Anchors, constructing semantic routes, or requiring additional manual prompts.

We evaluate RouteSeg on three 2D medical image segmentation datasets with different imaging characteristics and lesion types, namely Kvasir-SEG, BUSI, and ISIC2018, under a unified annotation budget of approximately \(1\%\). We compare RouteSeg with representative semi-supervised medical image segmentation methods, SAM-assisted low-annotation segmentation methods, and direct cross-image propagation strategies, so as to examine the role of \emph{cross-image supervision routing} in expanding limited supervision. In addition, we conduct ablation and analysis on key components such as Anchor selection, the introduction of Bridges, and semantic route construction, in order to study how different routing designs affect pseudo-supervision quality and final segmentation performance.

The main contributions of this paper are summarized as follows:

\begin{itemize}
    \item We propose RouteSeg, a cross-image supervision routing framework for semi-supervised medical image segmentation under extremely low annotation budgets. By viewing the unlabeled training set as a latent space formed by semantic relationships among samples, RouteSeg enables limited supervision to be gradually propagated through semantic routes.
    
    \item We propose Coverage-Aware Anchor Selection and an Anchor--Bridge--Target semantic routing mechanism. The former selects representative labeled samples as Anchors under a fixed annotation budget, while the latter retrieves semantically related unlabeled Bridges and performs progressive cross-image propagation based on SAM3, decomposing difficult direct propagation into multiple local steps.
    
    \item We propose reliable route-derived pseudo-supervision learning and a generalist--specialist collaboration strategy, which selects reliable pseudo-labels from mask candidates produced by different semantic routes and further improves the generalist's propagation results using task-specific knowledge learned by the specialist.
\end{itemize}


\section{Related Work}

\subsection{Semi-Supervised Medical Image Segmentation}

Semi-supervised medical image segmentation (SSMIS) aims to reduce annotation costs by jointly exploiting limited labeled data and abundant unlabeled images. Existing approaches mainly rely on pseudo-label learning, consistency regularization, and teacher--student paradigms. Cross Pseudo Supervision, Cross Teaching, and MC-Net+ improve unlabeled data utilization by enforcing prediction consistency among different networks or perturbations. Recent methods further enhance interactions between labeled and unlabeled samples. For example, BCP introduces bidirectional Copy--Paste augmentation, while ABD and DyCON improve robustness against pseudo-label uncertainty through adaptive perturbation and uncertainty-aware consistency learning.
Although these methods have achieved promising performance, unlabeled images are generally treated as independent training instances that require reliable predictions or regularization. How to exploit the internal semantic structures among large-scale unlabeled datasets for organizing supervision transfer remains less explored.


\subsection{Foundation Model-Assisted Label-Efficient Medical Segmentation}

The emergence of Segment Anything Model (SAM) provides new opportunities for label-efficient medical image segmentation. However, the domain gap between natural and medical images limits its direct application, motivating recent studies to combine SAM with semi-supervised learning.
SemiSAM and SemiSAM+ explore SAM-assisted semi-supervised segmentation by using task-specific models to generate prompts and leveraging SAM predictions as additional supervision. SemiSAM+ further introduces a Generalist--Specialist collaborative learning paradigm, where foundation models and task-specific networks provide complementary guidance. Other approaches, such as SAM-Match and CPAC-SAM, improve pseudo-label generation through SAM-based prompt interaction and adaptive sampling. Recent works including SC-SAM and CPPS-SAM further investigate bidirectional knowledge transfer between SAM and conventional segmentation networks.
These studies mainly focus on improving pseudo supervision quality or enhancing interactions between foundation models and specialist models. However, they usually consider unlabeled images individually and do not explicitly investigate how unlabeled samples can provide intermediate semantic transitions for supervision propagation.


\subsection{Cross-Image Semantic Modeling and Propagation}

Several studies have explored relationships among medical images to improve semi-supervised representation learning. ACTION utilizes anatomical-aware contrastive distillation to capture global semantic relationships and local anatomical characteristics. DACL constructs density-aware neighbor graphs to improve feature compactness, while GraphCL exploits graph structures among samples for representation enhancement. These methods demonstrate that medical images contain valuable cross-sample information, but such relationships are mainly used to optimize feature representations rather than determine how segmentation supervision should propagate across images.
Another line of research investigates cross-image knowledge transfer. OP-SAM generates semantic priors from support-query correlations and transfers knowledge from a single annotated image to target images. YoloSeg utilizes SAM2 propagation to expand one labeled image into multiple pseudo-labeled samples and reduces propagation noise through multi-view learning. However, these methods mainly focus on direct transfer between labeled sources and target images:
\begin{equation}
A \rightarrow T .
\end{equation}
In contrast, RouteSAM investigates whether abundant unlabeled images can serve as semantic bridges between a labeled Anchor and a Target. By constructing:
\begin{equation}
A \rightarrow B_1 \rightarrow B_2 \rightarrow \cdots \rightarrow T ,
\end{equation}
RouteSAM transforms unlabeled data from passive prediction targets into active components of supervision propagation, enabling gradual semantic transitions instead of direct supervision transfer.



\section{}\label{}

% To print the credit authorship contribution details
\printcredits

%% Loading bibliography style file
%\bibliographystyle{model1-num-names}
\bibliographystyle{cas-model2-names}

% Loading bibliography database
\bibliography{cas-refs.bib}

% Biography
%\bio{}
% Here goes the biography details.
%\endbio

%\bio{pic1}
% Here goes the biography details.
%\endbio

\end{document}

